\documentclass[10pt,a4paper]{article}
\usepackage[top=0.52in,bottom=0.48in,left=0.58in,right=0.58in]{geometry}
\usepackage[T1]{fontenc}
\usepackage[utf8]{inputenc}
\usepackage{lmodern}
\usepackage{xcolor}
\usepackage{titlesec}
\usepackage{enumitem}
\usepackage{tabularx}
\usepackage{hyperref}
\usepackage{fancyhdr}
\usepackage{microtype}

\definecolor{ink}{HTML}{122033}
\definecolor{muted}{HTML}{53657A}
\definecolor{blue}{HTML}{176FCA}
\definecolor{line}{HTML}{D8E1EB}

\hypersetup{colorlinks=true,urlcolor=blue,linkcolor=blue,pdftitle={Samitha Bandara - Curriculum Vitae},pdfauthor={Samitha Bandara}}
\pagestyle{fancy}
\fancyhf{}
\renewcommand{\headrulewidth}{0pt}
\fancyfoot[L]{\scriptsize\color{muted}Samitha Bandara}
\fancyfoot[R]{\scriptsize\color{muted}\thepage}
\titleformat{\section}{\large\bfseries\color{ink}}{}{0pt}{}[{\color{line}\titlerule[0.8pt]}]
\titlespacing*{\section}{0pt}{9pt}{5pt}
\setlist[itemize]{leftmargin=14pt,itemsep=2pt,topsep=3pt,parsep=0pt,label={\color{blue}\small\textbullet}}
\setlength{\parindent}{0pt}
\setlength{\parskip}{0pt}
\renewcommand{\arraystretch}{1.15}

\newcommand{\entry}[4]{%
  \textbf{\color{ink}#1}\hfill{\small\color{muted}#2}\\[-1pt]
  {\small\color{muted}#3}\hfill{\small\color{muted}#4}\par
}
\newcommand{\project}[3]{%
  \textbf{\color{ink}#1}\if\relax\detokenize{#2}\relax\else\hspace{5pt}{\small\href{#2}{GitHub}}\fi\\[-1pt]
  {\small\color{muted}#3}\par
}

\begin{document}
\begin{center}
  {\fontsize{26}{29}\selectfont\bfseries\color{ink} Samitha Bandara}\\[4pt]
  {\normalsize\color{muted}Final-Year Engineering Undergraduate \textbar{} AI Researcher}\\[7pt]
  {\small Sri Lanka \quad\textcolor{line}{\textbar}\quad
    \href{mailto:samithasahanssb@gmail.com}{samithasahanssb@gmail.com} \quad\textcolor{line}{\textbar}\quad
    \href{https://github.com/samitha278}{GitHub} \quad\textcolor{line}{\textbar}\quad
    \href{https://linkedin.com/in/samitha-bandara-52211a254}{LinkedIn} \quad\textcolor{line}{\textbar}\quad
    \href{https://samitha278.github.io}{Portfolio}}
\end{center}
\vspace{2pt}
\color{ink}

\section{Research Profile}
{\small Final-year Electronic and Telecommunication Engineering undergraduate researching efficient multimodal AI, visual representation reuse, 3D/4D perception, and learning-based robotic control. Interested in resource-aware models that perceive, reason, and act effectively, with the long-term goal of pursuing a PhD and an academic research career.}

\section{Research Experience}
\entry{Research Assistant}{}{Singapore Management University, School of Computing and Information Systems}{Singapore}
\begin{itemize}
  \item Investigated efficient vision-language models for repeated and multi-turn visual interactions through visual information reuse and compression.
  \item Explored Discrete Cosine Transform (DCT) based spectral compression, layer-wise reconstruction adapters, learned selection mechanisms, and visual attention or KV-cache reuse.
  \item Conducted experiments with multimodal models and visual question-answering benchmarks, focusing on representation efficiency and end-to-end inference cost.
  \item Contributed to two manuscripts under review at a leading international AI conference: first author on one and co-author on another.
\end{itemize}
{\small\color{muted}Research supported through the SMART Mens, Manus, and Machina (M3S) programme. Supervised by Prof. Archan Misra, with close guidance from Dr. Dulanga Weerakoon. Worked closely with Prasenjit Karmakar, PhD candidate at IIT Kharagpur.}

\section{Current Research Projects}
\project{Representation Reuse for Point-Cloud Vision-Language Models}{}{Exploring transformation-aware reuse of visual representations across dynamic point-cloud scenes to reduce redundant processing as viewpoints and scene geometry evolve.}
\vspace{5pt}
\project{Learning-Based Stabilization for an Aerial Manipulator}{}{Final-year project developing reinforcement-learning-based stabilization for a simulated aerial platform with a robotic arm, using PPO and Gazebo. Research collaboration with Prof. Rajat Talak, National University of Singapore.}

\section{Education}
\entry{BSc Engineering Honours in Electronic and Telecommunication Engineering}{Expected 31 May 2027}{University of Moratuwa, Sri Lanka}{CGPA: 3.79 / 4.00}
\begin{itemize}
  \item CGPA calculated across 110 completed GPA credits; Dean's List for Semesters 1, 2, and 3.
  \item Relevant coursework: Pattern Recognition, Image Processing and Machine Vision, Data Structures and Algorithms, Linear Algebra, Applied Statistics, Signals and Systems, Electronic Control Systems, and Communication Networks.
\end{itemize}
\vspace{3pt}
\entry{GCE Advanced Level - Physical Science}{2022}{Nalanda College, Colombo, Sri Lanka}{3 A grades}
{\small\color{muted}Z-score: 2.3862 \textbar{} Island rank: 278 \textbar{} Combined Mathematics, Physics, and Chemistry.}

\section{Research Interests}
{\small Efficient deep learning \textbar{} Vision-language and multimodal models \textbar{} 3D/4D perception \textbar{} Robot perception and intelligence \textbar{} Reinforcement learning \textbar{} Computer vision \textbar{} World models \textbar{} Mathematics and signal processing for machine learning}

\section{Selected Technical Projects}
\project{Transformer Optimization}{https://github.com/samitha278/transformer-optim}{Explored LoRA/PEFT adapters, 4-bit and 8-bit quantization, KV caching, and practical efficiency techniques for transformer models.}
\vspace{6pt}
\project{CoreLlama - Llama 2 from Scratch}{https://github.com/samitha278/CoreLlama}{Implemented RMSNorm, Rotary Positional Embeddings, Grouped Query Attention, KV caching, and SwiGLU directly in PyTorch.}
\vspace{6pt}
\project{VLM-gamma - Multimodal Vision-Language Model}{https://github.com/samitha278/VLM-gamma}{Built a multimodal model from scratch using a Vision Transformer encoder, transformer decoder, and Grouped Query Attention.}
\vspace{6pt}
\project{gpt2-lite - GPT-2 Reproduction}{https://github.com/samitha278/gpt2-lite}{Reproduced GPT-2 with a custom training pipeline covering mixed precision, compiled execution, optimized attention, gradient accumulation, scheduling, validation, and BPE tokenization.}
\vspace{6pt}
\project{nanoViT - Vision Transformer from Scratch}{https://github.com/samitha278/nanoViT}{Implemented patch embeddings, multi-head self-attention, transformer blocks, and image-classification pipelines.}
\vspace{6pt}
\project{MLP Character-Level Language Model}{https://github.com/samitha278/multilayer-perceptron-char-lm}{Implemented a character-level neural language model with backpropagation, batch normalization, Kaiming initialization, training, and evaluation.}

\section{Academic Engineering Projects}
\project{Battery Profiler - STM32 Embedded System}{https://github.com/samitha278/Battery_Profiler}{Designed a PCB with voltage and current sensing and a MOSFET load; implemented STM32 firmware for real-time acquisition and battery characterization.}
\vspace{5pt}
\project{Analog Five-Band Equalizer}{https://github.com/samitha278/5-Band-Equalizer}{Designed frequency-selective analog filter stages, gain controls, and a tested PCB prototype for real-time audio equalization.}

\section{Technical Skills}
\begin{tabularx}{\textwidth}{@{}>{\small\bfseries\color{ink}}p{0.22\textwidth} >{\small\color{muted}}X@{}}
Programming & Python, C, C++, MATLAB \\
Machine Learning & PyTorch, Hugging Face, TensorFlow, NumPy, Pandas, OpenCV \\
Research Topics & Vision-language models, transformers, computer vision, reinforcement learning, 3D/4D perception, model compression \\
Systems & Linux, Slurm, Docker, GPU training and optimization, Git \\
Robotics & Gazebo, PPO, simulation-based control, robotic perception \\
Embedded Engineering & STM32, embedded C, PCB design, sensors, data acquisition \\
Documentation & LaTeX, technical writing, Microsoft Office \\
Languages & English, Sinhala \\
\end{tabularx}

\section{Leadership and Activities}
\begin{itemize}
  \item Group leader for the final-year aerial-manipulator research project; responsible primarily for the reinforcement-learning component.
  \item Prefect, Nalanda College (2017--2019); participated in athletics, football, the Interact Club, and the Photographic and Art Society.
  \item Member of the University of Moratuwa Nature Team.
\end{itemize}
\end{document}
